\documentclass{article}
\usepackage{pseudocode}

\begin{document}

\section{compression}

\begin{algorithme}
    \fonction{compresser}{fic : FichierTexte}{fichierBinaire}
    {}
    {file : FileDePriorite;temp, table : TableDeCodage;  arbreHuffman : ArbreHuffman; fc : FichierBinaire}

    {
    \affecter{file}{creerFileDePriorite(fic)}
        creerFileDePriorite{fic}
    \affecter{table}{creerTableDeCodage(FileDePriorite.obtenirElement(file))}
    \tantque{non(EOF(fic))}{
        \affecter{temp}{table}
        ecrireOctet(fc,encoder(obtenirCodeBinaire(temp,lireCaractere(fic))))
    }

    }
    

    
    
    
\end{algorithme}

\section{creerFileDePriorite}
\begin{algorithme}
    \fonction{ }{fic :  FichierHuffman}{Arbre}
    {}
    {file : FileDePriorite; ensC : Ensemble de \caractere; ensS : Ensemble de Statistiques}
    {\affecter{ensC}{obtenirCaracteres(fic)}
    \affecter{ensS}{creerStatistique(ensC)}
    \affecter{file}{filedePriorite()}
    \pourChaque{Stats}{ensS}
    {enfiler(file, arbrehuffman(stats))}
    
    \tantque{FileDePriorite.obtenirLongueur(file)$>$1}  %OucreerFileDePriorite faire avec longueur >1 de FileDePrio
        {nouvelleFile(file)}
    \retourner(file)}
\end{algorithme}

\section{creerTableDeCodage}

\begin{algorithme}
    \fonction{ creerTableDeCodage}{arbre : ArbreHuffman}{TableDeCodage}{}{}
    {\retourner{obtenirTableDeCodage(arbre)}}
\end{algorithme}

\section{creerStatistique}

\begin{algorithme}
    \fonction{ creerStatistique}{fic : FichierTexte;ensC : Ensemble de \caractere}{Ensemble de Statistiques}
    {}
    {ensStats : Ensemble de Statistiques, nbOcc : \naturel}   				
    {\pourChaque{car}{ensC}
        \tantque{non(EOF(fic)){
        lire(fic,c)
        \sialors{c=car}{
            \affecter{nbOcc}{nbOcc + 1}
        Ensemble.ajouter(ensStats,statistique(car, nbOcc))}
    \retourner{ensStats}}
        }}
\end{algorithme}

\end{document}